\documentclass[]{../../../../tex/classes/homework}
\usepackage{../../../../tex/packages/hebrewSupport}


\author{שחר פרץ}
\title{חברה  דמוקרטית, תרבות וחינוך $\sim$ \textit{מטלה לסיום שיעור 7}}
\begin{document}
	\maketitle
	
	\section{}
	\textbf{שאלה: }מהם לדעתכם היתרונות והמגבלות בהגדרת סטנדרטיים פוליטיים ללימודי אזרחות, כמתואר בקטע שקראתם?
	
	\textbf{תשובה: }
	
	לדעתי יש מספר יתרונות להצבת סטנדרטיים פוליטיים ללימודי אזרחות, שאינם מצויינים בקטע. כמו שכאשר דרגים פוליטיים יכולים לממש את מדיניותם בנושאים אחרים מאזרחות על מנת לקדם מדיניות כוללת יותר (לדוגמה: שיפור לימודי הנדסה לצד תמיכה בחברת הנדסה אזרחיות, וכו') כך גם הצבת סטנדרטיים פוליטיים על לימודי האזרחות יכולה להשתלב לכדי תוכנית משמעותית יותר. לדוגמה, אם המחוקק סבור שישנה בעיה בלכידות החברה, הוא יכול לקדם מדיניות רפובליקאית יותר, לצד הפעלת תוכניות אזרחיות המקדמות זהות משותפת; ואם הוא חושב שישנה פגיעה באוטונומיה של קבוצות מסוימות, הוא יכול לקדם מדיניות ליברלית יותר, לצד קידום המדיניות מחוץ לבתי הספר. כך פתרון לבעיה אזרחית מוצע באופן יעיל יותר, שכן יש שילוב והתאמה בין החלטת הפוליטיקאי ללימודי אזרחות לצד פעולות שבסמכותו לעשות מחוץ לבתי הספר כדי לקדם את אותן המטרות. 
	
	עם זאת, ישנם גם חסרונות רבים, שעולים על היתרונות. במאה ה־21 (בין היתר הודות לפוליטיקת זהויות) מרבית האנשים מחולקים (בארה''ב, שהוא נושא הדיון של הקטע שקראנו) מחולקים לשני גושים – רפובליקאים (שלכאורה תומכים במדיניות רפובליקאית) ודמוקרטים (שלכאורה תומכים במדיניות ליברלית). לכן, לא רק שההתמקדות אך ורק ב\textit{שתי} גישות (רפובליקאית וליברלית) תצמצם את הדיון (כמתואר בקטע), אלא גם שכשמתאפשר לגורמים פוליטיים להשפיע על סטנדרטי למידה, הם לרוב יחילו סטנדרטים על לימודי האזרחות המציגים גישה \textit{אחת בלבד} (חינוך מוטה גישה – גם מצוין בקטע). כך, לדוגמה, אדם שלמד אזרחות במסגרת ליברלית יהיה פחות מודע לחשיבות השתתפותו שלו במנגנוני המדינה, ואדם שלמד במערכת רפובליקאית עלול לצאת בגישה אתנו־לאומית. כמתואר בטקסט, הניסיון לעודד את התלמידים לנקוט בגישה ספציפית (רפובליקאית או ליברלית) עלולה לגרום לסילוף או השמטת עובדות היסטוריות ויצירת תודעה שגויה של ההיסטוריה (לדוגמה, הלימודים בישראל שמתעלמים מהחשיבות של תנועות נוער לא־ציוניות בשואה, הלימודים בצרפת שממעיטים משלטון הטרור במהפכה הצרפתית, הזנחה של אירועים בין לאומיים בתוכנית הלימודים בבריטניה, ודוגמאות מהקטע. עשיתי על זה מתישהו עבודה לבית הספר). העובדה שלימודי ההיסטוריה והאזרחות בארה''ב לרוב שניהם מחוברים ללמודי ''Social Studies``, מקצוע יחיד, לא עוזרת, ומעודדת שילוב נקודת מבט לאומית (רגשית במידת מה) על ההיסטוריה העובדתית. 
	
	סך הכל, הגדרת סטנדרטים פוליטיים על לימודי האזרחות יכולה להועיל לפתרון בעיות אזרחיות באופן יעיל באמצעות שילוב של מאמצים לפתרון בעיות אלו מחוץ לבתי הספר, באמצעות תיאום של גורמים פוליטיים. אך מצד שני, הגדרת הסטנדרטיים האלו מגבילה את הדיון לשתי גישות בלבד, או כמו שלרוב קורה (לאור הפילוג לשני מחנות) לגישה אחת בלבד. היא גם עלולה לגרור הבנה שגויה של העובדות ההיסטוריות. 
	
	\section{}
	\textbf{שאלה: }מהם קווי הדמיון והשוני בין הגישה הרפובליקאית והליברלית לאזרחות?
	
	\textbf{תשובה: }
	
	אומנם הגישות הרפובליקאיות והליברלית שונות אחת מהשנייה, אך הן גם חולקות קווי דמיון. בין היתר, שתי הגישות מדגישות את החשיבות של שלטון החוק ושל קבלת החלטות בתוך המנגנון הדמוקרטי, בניגוד לתפיסות אחרות שמייחסות פחות חשיבות למנגנונים מדינתיים (או מתנגדות לקיומם של כאלו, לדוגמה, אנרכיזם). שתי הגישות מכירות במוסד המדינה כמוסד המחוקק שאינו כפוף לסמכות של מוסדות גדולים יותר (בניגוד לגישה העל־לאומית שתומכת בדלגציה של סמכויות לגופים בין־לאומיים) ואינן מדירות קבוצות ספציפיות על בסיס מין, דת, תרבות וכו' (באפילו מקרה הרפובליקאי, הדגש הוא על כך שהאזרח רואה את עצמו כשייך למדינה, לא לאיזו קבוצה אתנית האזרח שייך). שתי הגישות תומכות בהבנה של המצב הפוליטי, אם כי הגישה הליברלית מתמקדת בדיון עליו בעוד הגישה הרפובליקאי דוגלת בהנחלת ידע. 
	
	\npage
	
	בין שתי הגישות יש גם הבדלים רבים. לדוגמה, הגישה הרפובליקאית מציבה את האינטרסים של הקהילה במרכז, לפני האינטרסים של הפרט. בעולם בו מוקד זהותי חזק הוא הלאום, במקרים רבים מדובר על הצבת המדינה לפני האדם. מבחינתה האזרחות כוללת תמיכה בממשל ופיתוח זהות התואמת זאת. הגישה הליברלית שמה דגש על החופש של כל אדם ללא תלות לאיזה לאום או מוקד זהות הוא שייך אליו, ולדידה האזרחות כוללת מאפיינים של חופש זכויות שהממשל מספק לפרט (לדוגמה, באמצעות הגנה על זכויותיו, באמצעים כמו מערכות אכיפת החוק). הגישה הליברלית מעודדת השתתפות פעילה בהליך הדמוקרטי והשפעה עליו. 
	
	לסיום, יש קווי דימיון רבים בין הגישה הרפובליקאית לגישה הליברלית, באופן שבו הן מכירות במוסד המדינה ובסמכויותיו. למרות זאת, הגישה הרפובליקאית מציבה את הממסד והקהילה כבעלי חשיבות עליונה, לאומת הגישה הליברלית ששמה דגש על החופש של הפרט בתוך החברה, ומכאן גם נגזרת גישתן לאזרחות. 
	
	
	
	
	\ndoc
\end{document}